% ============================================================
%  HALLAZGOS  (HAL-001 – HAL-007)
% ============================================================
\clearpage
\section{Hallazgos Encontrados}

\subsection{HAL-001: Producto inactivo sigue visible en módulo de Ventas}
\hallazgo
  {HAL-001}
  {Producto inactivo visible en Ventas}
  {\sevAlta}
  {Productos \arr{} Ventas}
  {Bug funcional}
  {CP-M006}
  {Al marcar un producto como inactivo en el módulo de Productos, dicho producto
   continúa apareciendo disponible para seleccionar al crear una nueva venta.}
  {Se pueden registrar ventas de productos dados de baja intencionalmente, generando
   inconsistencias en el inventario y los reportes.}
  {Agregar filtro \texttt{WHERE PRO\_Estado = 'A'} (activo) en la consulta que carga
   productos disponibles para venta.}

\subsection{HAL-002: Buscador de productos no filtra por código}
\hallazgo
  {HAL-002}
  {Buscador de productos no filtra por código}
  {\sevMedia}
  {Productos}
  {Funcionalidad incompleta}
  {CP-A003}
  {El buscador del módulo de Productos solo filtra por nombre. Al ingresar un código
   de producto (ej: ``ARE-001''), no retorna ningún resultado aunque el producto exista
   con ese código.}
  {Los usuarios que conozcan el código del producto pero no el nombre exacto no pueden
   encontrarlo, reduciendo la eficiencia operativa.}
  {Incluir el campo \texttt{PRO\_Codigo} en el filtro de búsqueda, de forma que busque
   por nombre \textbf{O} código.}

\subsection{HAL-003: No existe historial de movimientos de inventario}
\hallazgo
  {HAL-003}
  {No existe historial de movimientos de inventario}
  {\sevAlta}
  {Inventario}
  {Funcionalidad faltante}
  {CP-A004}
  {El módulo de Inventario permite registrar entradas y salidas de mercadería, pero no
   existe ninguna opción para consultar el historial de movimientos registrados
   (auditoría de entradas, salidas y ajustes).}
  {No es posible hacer trazabilidad de los movimientos de inventario. No se puede
   verificar quién realizó qué movimiento ni cuándo, lo cual es crítico para auditorías
   y control interno.}
  {Agregar una sección ``Historial de movimientos'' en el módulo de Inventario que liste
   los movimientos con fecha, tipo, cantidad, producto y usuario responsable.}

\subsection{HAL-004: No hay indicador visual de stock bajo en módulo Inventario}
\hallazgo
  {HAL-004}
  {Sin indicador visual de stock bajo en Inventario}
  {\sevMedia}
  {Inventario}
  {Mejora de usabilidad}
  {CP-A009}
  {El módulo de Inventario muestra el stock de todos los productos, pero no diferencia
   visualmente los productos con stock bajo o agotado. El Dashboard sí muestra el
   contador ``Stock Bajo'', pero al entrar al módulo de Inventario no hay colores,
   íconos ni etiquetas que identifiquen productos críticos.}
  {El usuario debe revisar manualmente cada producto para identificar cuáles necesitan
   reabastecimiento, aumentando el riesgo de quedarse sin stock de productos importantes.}
  {Agregar indicadores visuales (color rojo/amarillo, ícono de alerta) para productos
   con stock igual a 0 o por debajo de un umbral definido.}

\subsection{HAL-005: Campo teléfono en Emprendedoras acepta letras}
\hallazgo
  {HAL-005}
  {Campo teléfono en Emprendedoras acepta letras}
  {\sevBaja}
  {Emprendedoras}
  {Validación faltante}
  {CP-A008}
  {El campo de teléfono en el formulario de Emprendedoras acepta cualquier carácter,
   incluyendo letras. No existe validación de formato numérico.}
  {Se pueden guardar números de teléfono inválidos (ej: ``abc123''), generando datos
   de contacto incorrectos e inutilizables.}
  {Agregar validación en el campo teléfono para aceptar solo dígitos numéricos, con
   un formato esperado de 8 dígitos (estándar Costa Rica).}

\subsection{HAL-006: US-14 (Gestión por local) no está implementada}
\hallazgo
  {HAL-006}
  {Gestión por local no implementada (US-14)}
  {\sevMedia}
  {---}
  {Funcionalidad no implementada}
  {CP-A014}
  {La User Story US-14 requiere gestionar información de diferentes locales o puntos de
   venta. No existe ningún módulo, sección ni campo relacionado con locales en la
   interfaz de la aplicación.}
  {Si la asociación necesita expandirse a múltiples locales, el sistema no tiene soporte
   para esto.}
  {Implementar el módulo de gestión de locales según los requerimientos de la US-14,
   incluyendo creación, edición y asignación de inventario por local.}

\subsection{HAL-007: Botón ``Destacar producto'' no funciona}
\hallazgo
  {HAL-007}
  {Botón ``Destacar producto'' no funciona}
  {\sevAlta}
  {Productos}
  {Bug funcional}
  {CP-A016}
  {En el módulo de Productos existe un botón para marcar un producto como
   ``destacado'' para mostrarlo en la página web. Al presionar el botón, no ocurre
   ninguna acción visible: no cambia de estado, no muestra confirmación ni error.}
  {La User Story US-16 no puede cumplirse. Los administradores no pueden promover
   productos en el sitio web desde la aplicación.}
  {Revisar el handler del botón de destacado en el componente de productos y verificar
   que el IPC correspondiente esté registrado y funcionando correctamente.}

\subsection{Tabla Resumen de Hallazgos}

\begin{tabularx}{\linewidth}{|c|c|X|l|c|}
\hline
\rowcolor{grisheader}
\textbf{ID} & \textbf{Severidad} & \textbf{Módulo} & \textbf{Tipo} & \textbf{Estado}\\
\hline
HAL-001 & \sevAlta  & Productos \arr{} Ventas & Bug funcional            & Confirmado\\
\hline
HAL-002 & \sevMedia & Productos               & Funcionalidad incompleta  & Confirmado\\
\hline
HAL-003 & \sevAlta  & Inventario              & Funcionalidad faltante    & Confirmado\\
\hline
HAL-004 & \sevMedia & Inventario              & Usabilidad                & Confirmado\\
\hline
HAL-005 & \sevBaja  & Emprendedoras           & Validación faltante       & Confirmado\\
\hline
HAL-006 & \sevMedia & ---                     & No implementado           & Confirmado\\
\hline
HAL-007 & \sevAlta  & Productos               & Bug funcional             & Confirmado\\
\hline
\end{tabularx}
